% !TeX root = ../main.tex
 
\chapter{Algorithmen}

\section{Motivation}

\dfn{Elementare Addition} {
  Addiere drei 1-stellige Zahlen und erhalten ein 2-stelliges Ergebnis
}

\nt{Zur naiven Addition zweier $n$-stelligen Zahlen benötigt man $n$ elementare Additionen. Der schnellste bekannte Additionsalgorithmus ist $\log_2 n + \log_2 \log_2 n + \log_2 \log_2 \log_2 n + 5.6$}

\dfn{Elementare Multiplikation}{
  Multipliziere zwei 1-stellige Zahlen und erhalten ein 2-stelliges Ergebnis
}

\nt{Man weiß nicht wie schnell man Multiplizieren kann.}
\ex{LGS}{
  \begin{align*}
    10^{-20}x + 2y &= 1 \\
    \underline{10^{-20}x + 10^{-20}y &= 10^{-20}} \\
    (2-10^{-20})y &= 1-10^{-20} \\
    \text{Taschenrechner sagt:} \\
    \implies y &\approx 0.5 \\
    \text{Setze in die erste Gleichung:} \\
    \implies x &= 0 \\
    \text{Setze in die zweite Gleichung:} \\
    \implies x &= 0.5
  \end{align*}
  Korrekt wäre
  \begin{align*}
    x &= \frac{1}{2-10^{-20}} = 0.500\dots \\
    y &= \frac{1-10^{-20}}{2-10^{-20}} = 0.499\dots
  \end{align*}
}

\section{2}

\dfn{Kartesisches Produkt}{
  $$
  A, B := \text{Mengen}
  $$
  Dann ist deren \textbf{Kartesisches Produkt} definiert durch $A \times B := \{(a,b)\} \vert a \in A, b \in B$. \\
  Eine \textbf{Relation} auf $(A, B)$ ist eine Teilmenge von $A \times B$. \\
  Schreibweise $a R b$
}
\ex{}{
  $$
    A = \{1, 3, 5\}, B=\{2, 5\}\\
    A \times B = alle Permutationen
    R = Teilmenge
  $$
}
